\documentclass[10pt]{article}

% ---------- Document Identity (update these only) ----------
\newcommand{\DocStage}{}
\newcommand{\DocTitle}{Your Road to Tier One SOF}
\newcommand{\ProgramName}{182-Week Civilian-to-Selection Readiness Pipeline}
\newcommand{\ProgramSubtitle}{}

% ---------- Page ----------
\usepackage[legalpaper,left=0.5in,right=0.5in,top=0.6in,bottom=0.45in]{geometry}

% ---------- Typography ----------
\usepackage[T1]{fontenc}
\usepackage[utf8]{inputenc}
\usepackage{libertinus}
\usepackage{microtype}
\usepackage{parskip}
\usepackage{ragged2e}
\usepackage{textcomp}
\AtBeginDocument{\color{black}}
\AtBeginDocument{\pagecolor{white}}
\AtBeginDocument{\justifying}

% ---------- Landscape pages ----------
\usepackage{pdflscape}

% ---------- Color ----------
\usepackage[table]{xcolor}
\definecolor{StageBlue}{HTML}{1F4E79}
\definecolor{StageBlueDark}{HTML}{163A5A}
\definecolor{LightGray}{HTML}{F2F4F7}
\definecolor{MidGray}{HTML}{D9DEE5}
\definecolor{RuleGray}{HTML}{C7CED8}
\definecolor{HeaderInk}{HTML}{000000}
\definecolor{Ink}{HTML}{000000}

% ---------- Utility ----------
\usepackage{enumitem}
\sloppy
\setlength{\emergencystretch}{2em}

% ---------- Boxes / UI ----------
\usepackage[most]{tcolorbox}

\tcbset{
  charterTitle/.style={
    enhanced,
    colback=StageBlue!4,
    colframe=StageBlue,
    boxrule=1.25pt,
    arc=3pt,
    left=16pt,right=16pt,top=14pt,bottom=14pt,
    borderline north={3.2pt}{0pt}{StageBlue},
    borderline west={4pt}{0pt}{StageBlue},
    borderline south={1.25pt}{0pt}{StageBlue},
    drop shadow={black!25!white},
  },
  charterRule/.style={
    enhanced,
    colback=LightGray,
    colframe=RuleGray,
    boxrule=0.85pt,
    arc=3pt,
    left=12pt,right=12pt,top=10pt,bottom=10pt,
    borderline west={3pt}{0pt}{StageBlue},
  },
  charterBadge/.style={
    enhanced,
    colback=StageBlue,
    coltext=white,
    colframe=StageBlueDark,
    boxrule=0.6pt,
    arc=2pt,
    left=6pt,right=6pt,top=3pt,bottom=3pt,
  }
}
\newtcbox{\Badge}{nobeforeafter,charterBadge}

% ---------- Lists ----------
\setlist[itemize]{leftmargin=1.5em, itemsep=0.25em, topsep=0.25em}
\setlist[enumerate]{leftmargin=1.8em, itemsep=0.25em, topsep=0.25em}

% ---------- TOC ----------
\usepackage{tocloft}
\setcounter{tocdepth}{3}
\setlength{\cftbeforesecskip}{3pt}
\setlength{\cftbeforesubsecskip}{2pt}
\setlength{\cftbeforesubsubsecskip}{0pt}
\renewcommand{\cftsecfont}{\bfseries}
\renewcommand{\cftsecpagefont}{\bfseries}
\renewcommand{\cftsubsecfont}{\normalfont}
\renewcommand{\cftsubsecpagefont}{\normalfont}
\renewcommand{\cftsubsubsecfont}{\normalfont}
\renewcommand{\cftsubsubsecpagefont}{\normalfont}
\setlength{\cftsecindent}{0pt}
\setlength{\cftsubsecindent}{1.25em}
\setlength{\cftsubsubsecindent}{2.8em}
\setlength{\cftsecnumwidth}{2.1em}
\setlength{\cftsubsecnumwidth}{2.8em}
\setlength{\cftsubsubsecnumwidth}{3.6em}

% Remove the built-in TOC heading so only the custom heading is shown
\makeatletter
\renewcommand{\tableofcontents}{\@starttoc{toc}}
\makeatother

% ---------- Header/Footer: page number only ----------
\usepackage{fancyhdr}
\pagestyle{fancy}
\fancyhf{}
\renewcommand{\headrulewidth}{0pt}
\renewcommand{\footrulewidth}{0pt}
\fancyfoot[C]{\thepage}

% ---------- Section formatting ----------
\usepackage{titlesec}

\titleformat{\section}
  {\Large\bfseries\color{Ink}}
  {\thesection}{0.6em}{}
  [\vspace{0.25em}\textcolor{StageBlue}{\rule{\linewidth}{0.8pt}}]

\titleformat{\subsection}
  {\large\bfseries\color{Ink}}
  {\thesubsection}{0.6em}{}

\titleformat{\subsubsection}
  {\normalsize\bfseries\color{Ink}}
  {\thesubsubsection}{0.6em}{}

\titlespacing*{\section}{0pt}{0.95em}{0.35em}
\titlespacing*{\subsection}{0pt}{0.65em}{0.25em}
\titlespacing*{\subsubsection}{0pt}{0.55em}{0.2em}

% ---------- Table engine ----------
\usepackage{tabularray}
\UseTblrLibrary{booktabs}

% ---------- tabularray: GLOBAL table treatment ----------
\SetTblrInner{
  cells = {fg=black, bg=white, valign=t},
}

% ---------- Header helpers ----------
\newcommand{\Hdr}[1]{\shortstack[c]{\strut #1\strut}}

% ---------- Lines ----------
\newcommand{\TitleLine}{\textcolor{StageBlue}{\rule{0.98\linewidth}{0.95pt}}}
\newcommand{\SubTitleLine}{\textcolor{RuleGray}{\rule{0.98\linewidth}{0.65pt}}}

% ---------- Continued on next page ----------
\newcommand{\ContinuedOnNextPageInline}{%
  \par
  \vfill
  \noindent\textcolor{RuleGray}{\rule{\linewidth}{1pt}}\vspace{-2pt}
  \noindent\hfill\textit{Continued on next page}\par\vspace{-10pt}
  \noindent\textcolor{RuleGray}{\rule{\linewidth}{1pt}}
  \clearpage
}

% ---------- PDF hyperlinks ----------
\usepackage[hidelinks]{hyperref}
\hypersetup{
  colorlinks=false,
  pdfborder={0 0 0},
  linktoc=all,
  pdftitle={\DocTitle},
  pdfauthor={\ProgramName}
}

% =========================================================
\begin{document}

\section*{7.13 — Nutrition Preparation and Hydration Systems}
\phantomsection
\addcontentsline{toc}{section}{7.13 — Nutrition Preparation and Hydration Systems}

Overview \textbf{ID}: 7.13\\
Overview Name: Nutrition Preparation and Hydration Systems\\
Version: v1.0\\
Governing Status: Draft — Non-Deployable\\
Scope Boundary Statement: Covers food-preparation tools, storage infrastructure, hydration container systems, field hydration logistics, and cleaning / drying support tools; does not provide macro prescriptions, diet programming, supplement dosing, or medical guidance.

\subsection*{7.13.1 Purpose \& Scope}
\phantomsection
\addcontentsline{toc}{subsection}{7.13.1 Purpose \& Scope}

Category coverage:

\begin{itemize}
\item Food preparation tools and storage infrastructure (containers, basic staging, carry posture).
\item Hydration container systems and field hydration logistics (bottles and reservoirs with cleaning/drying tools).
\item Simple planning workflow supports that reduce friction and missed sessions (staging posture, repeatable cleaning-to-ready routine).
\end{itemize}

This overview crosswalks Stage 6 timing to the nutrition-preparation and hydration conditions required for safe execution and admissible evidence.

Exclusions:

\begin{itemize}
\item No macro/calorie prescriptions or dieting protocols (Stage 10 owns nutrition programming).
\item No supplement dosing guidance (cross-reference Stage 7.03–Stage 7.06 only).
\item No training programming content.
\end{itemize}

This overview does not provide programming, test protocols, calendars, macro prescriptions, diet protocols, supplement dosing, or medical guidance.

\subsection*{7.13.2 Governance And Safety Boundaries}
\phantomsection
\addcontentsline{toc}{subsection}{7.13.2 Governance And Safety Boundaries}

\begin{itemize}
\item Stage 0 boundary alignment: stop rules, safety override doctrine, conflict-resolution hierarchy, and evidence-admissibility controls apply to all 7.13 selections and substitutions.
\item Tools enable adherence; they do not replace nutrition strategy or medical care.
\item Food safety is a required control:
  \begin{itemize}
  \item clean hands and surfaces posture,
  \item separation of raw and ready-to-eat posture,
  \item rapid cooling and safe storage posture,
  \item discard-when-uncertain posture.
  \end{itemize}
  Minimal temperature conventions (general food-safety convention, not a medical protocol):
  \begin{itemize}
  \item cold foods \textbf{SHOULD} remain cold (refrigerated) and \textbf{SHOULD} be transported with cold packs when held for extended periods.
  \item hot foods \textbf{SHOULD} be kept hot or cooled promptly; do not leave perishable foods at room temperature for long periods.
  \end{itemize}
\item Hydration safety is a required control:
  \begin{itemize}
  \item avoid intentional dehydration practices,
  \item plan for safe access and daily cleaning,
  \item do not ignore dehydration symptoms; downshift and correct feasibility.
  \end{itemize}
\item Evidence inadmissibility: evidence produced under unsafe food-handling, contaminated storage, unresolved leakage, or hydration-system failure conditions is inadmissible until corrected and re-verified.
\item Change-control posture:
  \begin{itemize}
  \item one change at a time (primary bottle/reservoir, container system, workflow),
  \item no novelty near Deload+Test weeks and gate windows,
  \item log tool changes, start date, reason, and any issues (leaks, GI tolerance, cleaning failures).
  \end{itemize}
\end{itemize}

\subsection*{7.13.3 Tier Doctrine}
\phantomsection
\addcontentsline{toc}{subsection}{7.13.3 Tier Doctrine}

\begin{itemize}
\item Price band convention: \$ = minimal household-cost item; \$\$ = modest purpose-built item; \$\$\$ = expanded-capability or redundancy item.
\item N/A rule: use N/A only when a field is genuinely not applicable and omission does not obscure safety, maintenance, or phase-timing logic.
\item Substitution posture: substitutions are acceptable only when they preserve food safety, cleanability, leak control, and staging reliability; substitutions that increase contamination risk, spoilage risk, or workflow failure break intent.
\item Tier 0: household-only stopgaps with explicit safety caveats; minimal safe container reuse and basic bottle posture only when cleanable and intact.
\item Tier 1: minimum viable prep and hydration system that prevents missed sessions and reduces failure from logistics; includes cleaning tools (brush + drying rack).
\item Tier 2: expanded system improving repeatability and reducing spoilage risk; includes food scale (Tier 2) and hydration reservoir/bladder systems (Tier 2).
\item Tier 3: overbuilt redundancy and efficiency; convenience and sustainment bias without changing what is required.
\end{itemize}

\subsection*{7.13.4 Selection Criteria \& Fit Checks}
\phantomsection
\addcontentsline{toc}{subsection}{7.13.4 Selection Criteria \& Fit Checks}

Fit and safety checklist (operational):

Containers:

\begin{itemize}
\item Leakproofing and seal integrity: lids and gaskets must seal reliably.
\item Cleanability: wide-mouth access and smooth surfaces reduce residue retention.
\item Durability: resists cracking; stable in bag.
\item Dishwasher-safe posture is preferred where available; hand-wash compatibility must remain practical.
\end{itemize}

Bottles/reservoirs:

\begin{itemize}
\item Opening size supports cleaning and full drying.
\item Odor retention risk is minimized by clean-to-ready posture and full drying.
\item Transport compatibility: fits bags and does not leak under pressure changes.
\item Reservoir routing compatibility with Stage 7.07 (if used) and secure bite valve posture.
\end{itemize}

Storage/carry:

\begin{itemize}
\item Insulation and spill containment reduce spoilage risk.
\item Packability supports daily execution without friction.
\end{itemize}

Food scale:

\begin{itemize}
\item Stable surface use and wipe-down posture.
\item Battery replacement posture to prevent missing data events.
\end{itemize}

Fit checks (clean-to-ready workflow):

\begin{itemize}
\item The full loop must be feasible daily:
  \begin{itemize}
  \item prepare → carry → consume → rinse → scrub → dry fully → stage for next day.
  \end{itemize}
\item If the loop is not feasible, simplify rather than adding complexity.
\end{itemize}

\subsection*{7.13.5 Setup, Safety, and Anchoring}
\phantomsection
\addcontentsline{toc}{subsection}{7.13.5 Setup, Safety, and Anchoring}

Setup doctrine (no programming):

Food safety workflow posture (practical, non-medical):

\begin{itemize}
\item Clean hands and surfaces before prep.
\item Separate raw and ready-to-eat items; avoid cross-contamination.
\item Cool perishable items promptly; store cold items cold; discard-when-uncertain posture.
\item Do not “push through” questionable food; GI distress degrades training safety and validity.
\end{itemize}

Hydration setup posture:

\begin{itemize}
\item Daily rinse after use; periodic deep clean with brush tools; dry fully.
\item Reservoir-specific:
  \begin{itemize}
  \item fully drain and dry; do not store sealed while damp.
  \item replace valves/tubes when persistent odor or residue cannot be resolved.
  \end{itemize}
\item Gate-window stability:
  \begin{itemize}
  \item do not introduce a new bottle/reservoir system inside Deload+Test windows.
  \item if a system must change due to failure, log the change and treat decision-window comparability as higher variance; re-slot gate attempts if the change affects execution.
  \end{itemize}
\end{itemize}

Field logistics:

\begin{itemize}
\item Pre-pack hydration and food carry the night before protected windows to reduce day-of omission risk.
\item Verify leak checks before placing in ruck/pack.
\end{itemize}

\subsection*{7.13.6 Maintenance, Replacement, and Storage}
\phantomsection
\addcontentsline{toc}{subsection}{7.13.6 Maintenance, Replacement, and Storage}

Replacement cues:

\begin{itemize}
\item Persistent odor that persists after deep cleaning.
\item Visible residue that cannot be removed.
\item Cracked seals, leaking lids, warped container edges.
\item Reservoir: persistent mold or odor, degraded valves or tubes, or leaks.
\end{itemize}

Storage posture:

\begin{itemize}
\item Dry all components fully before closing or storing them in bins.
\item Store with airflow; avoid trapping moisture.
\item Keep containers staged clean-to-ready to prevent missed sessions.
\end{itemize}

Consumables boundary:

\begin{itemize}
\item Soap/sanitizers and hygiene consumables are Stage 7.20; cross-reference only.
\item 7.13 lists cleaning tools (brushes, drying rack) and non-chemical maintenance items (replacement lids/gaskets, ice packs).
\end{itemize}

\subsection*{7.13.7 Substitution Rules — Minimal Path Alternatives}
\phantomsection
\addcontentsline{toc}{subsection}{7.13.7 Substitution Rules — Minimal Path Alternatives}

\begin{itemize}
\item If no meal-prep system:
  \begin{itemize}
  \item use a simpler container system with fewer parts; reduce complexity rather than skipping meals.
  \end{itemize}
\item If no reservoir:
  \begin{itemize}
  \item bottle-based hydration remains acceptable until Stage 6 requires reservoir capability by Phase 06 (End).
  \end{itemize}
\item If a container fails:
  \begin{itemize}
  \item discard and replace rather than unsafe reuse.
  \end{itemize}
\item Validity and logging linkage (Stage 2.E posture, high level):
  \begin{itemize}
  \item log changes to primary hydration system as a tool change (tool description, start date, reason).
  \item document issues that impact execution (leakage, GI tolerance, cleaning failures).
  \item do not infer missing intake data; record \textbf{MISSING}/\textbf{UNKNOWN} where required by the program’s missing-data doctrine.
  \end{itemize}
\end{itemize}

\subsection*{7.13.8 Phase Integration Map — Required}
\phantomsection
\addcontentsline{toc}{subsection}{7.13.8 Phase Integration Map — Required}

\begingroup
\scriptsize
\setlength{\tabcolsep}{2pt}
\renewcommand{\arraystretch}{1.12}

\begin{longtblr}{
  width=\linewidth,
  rowhead=1,
  colspec = {
    Q[l,wd=1.05in]
    Q[l,wd=1.55in]
    X[l,2.75]
    X[l,3.65]
  },
  hlines,
  vlines,
  cells = {hsep=6pt, vsep=6pt, halign=l, valign=t},
  row{1} = {bg=StageBlue!15, fg=black, font=\bfseries, halign=l, valign=t},
  row{odd} = {bg=white},
  row{even} = {bg=LightGray},
}
\Hdr{Phase} &
\Hdr{Required Tier (from Stage 6)} &
\Hdr{What Changes / Becomes Mandatory} &
\Hdr{Notes} \\
Phase 00 &
Tier 1 &
Tier 1 hydration and basic food storage system is required by Phase 00 (End) &
Phase 00 (End): primary bottle + basic containers + brush/drying posture must exist to protect adherence and safety. No novelty in protected windows. \\
Phase 01 &
Tier 1 &
Maintain Tier 1; required clean-to-ready workflow consistency becomes operationally mandatory &
Phase 01 (End): missed hydration/food planning becomes a compliance failure driver; simplify workflow rather than skipping. \\
Phase 02 &
Tier 1 &
Maintain Tier 1; required stability supports higher consequence training load &
Phase 02 (End): do not add new container systems in protected windows; log changes. \\
Phase 03 &
Tier 1 &
Maintain Tier 1; required field carry discipline increases as ruck systems begin to appear &
Phase 03 (End): verify leak checks; cross-reference Stage 7.07 for integration when carried in ruck. \\
Phase 04 &
Tier 1 &
Maintain Tier 1; required cleaning/drying discipline increases with outdoor exposure &
Phase 04 (End): wet/dirty exposure increases spoilage and mold risk; drying posture is non-negotiable. \\
Phase 05 &
Tier 1 &
Maintain Tier 1; required carry reliability increases as endurance duration increases &
Phase 05 (End): longer sessions increase consequence of missing hydration/food logistics; do not improvise. \\
Phase 06 &
Tier 2 &
Tier 2 expansion is required by Phase 06 (End): reservoir system + food scale + expanded containers &
Phase 06 (End): reservoir/bladder becomes required; cross-reference Stage 7.07 routing / fit. Food scale supports prep repeatability only; no macros here. \\
Phase 07 &
Tier 2 &
Maintain Tier 2; required standardization reduces variance in hybrid density &
Phase 07 (End): stabilize reservoir system; log changes; no novelty near protected windows. \\
Phase 08 &
Tier 2 &
Maintain Tier 2; required robust cleaning posture protects from mold/odor failures &
Phase 08 (End): ruck/load phases punish hydration failures; clean-to-ready discipline is a safety control. \\
Phase 09 &
Tier 2 &
Maintain Tier 2; required throughput consistency reduces missed sessions &
Phase 09 (End): batch-prep capacity reduces friction; avoid food safety drift. \\
Phase 10 &
Tier 2 &
Maintain Tier 2; required field logistics become high consequence &
Phase 10 (End): protect against leakage and GI disruption; log tool changes. \\
Phase 11 &
Tier 2 &
Maintain Tier 2; required low variance supports recovery-sensitive phase &
Phase 11 (End): do not change container/reservoir systems in decision windows. \\
Phase 12 &
Tier 2 &
Maintain Tier 2; optional Tier 3 redundancy may reduce compromised windows &
Phase 12 (End): if redundancy is adopted, stabilize well before verification blocks; never introduce new systems inside protected windows. \\
Phase 13 &
Tier 2 &
Maintain Tier 2; consolidation requires stable systems and clean-to-ready execution &
Phase 13 (End): stability is the objective; avoid last-minute system changes. \\
\end{longtblr}

\endgroup

7.13.8 Phase Integration Map Notes:

\begin{itemize}
\item If any upstream phase artifact demands reservoir/bladder capability earlier than Phase 06 (End), requires Stage 8 review.
\item If any artifact attempts to turn 7.13 into macro or diet prescriptions, requires Stage 8 review (Stage 10 owns nutrition programming).
\end{itemize}

\subsection*{7.13.9 Risks, Contraindications, and Red Flags}
\phantomsection
\addcontentsline{toc}{subsection}{7.13.9 Risks, Contraindications, and Red Flags}

Non-medical risks:

\begin{itemize}
\item Food spoilage and GI disruption from poor handling:
  \begin{itemize}
  \item control: discard-when-uncertain posture; avoid long room-temperature holds for perishables; clean surfaces and separate raw/ready-to-eat.
  \end{itemize}
\item Mold/odor and contamination risk from poorly cleaned bottles/reservoirs:
  \begin{itemize}
  \item control: rinse daily, deep clean periodically, dry fully; never store sealed while damp.
  \end{itemize}
\item Leakage causing missed sessions or ruined gear:
  \begin{itemize}
  \item control: leakproofing and gasket discipline; staged spares; leak check before packing.
  \end{itemize}
\item Underhydration from missing containers:
  \begin{itemize}
  \item control: backup bottle posture; staging the night before high-consequence days.
  \end{itemize}
\item Red flags requiring conservative stop/escalation posture:
  \begin{itemize}
  \item repeated GI distress tied to food-handling failures,
  \item visible mold or persistent odor that cannot be resolved,
  \item repeated dehydration symptoms during sessions.
  \end{itemize}
\item Requires Stage 8 review triggers:
  \begin{itemize}
  \item any mismatch requiring Tier 1 or Tier 2 capability earlier than Stage 6 timing,
  \item any attempt to include macro prescriptions or diet protocols inside 7.13,
  \item any attempt to list hygiene consumables as Stage 7.13 consumables rather than cross-referencing Stage 7.20.
  \end{itemize}
\end{itemize}

\end{document}